\section{Introduction}

\begin{figure}[!t]
\centering
\includegraphics[width=\columnwidth]{figures/fig1-rho-comparison.pdf}
\caption{\looseness=-1 \method{} versus validation-feedback harness optimization.
Validation-feedback methods iterate against a labeled validation set, whereas \method{} optimizes from past trajectories in a single retrospective pass with no ground-truth labels.}
\vspace{-1.5em}
\label{fig:rho-comparison}
\end{figure}

\looseness=-1 A \textit{harness} enables an AI agent to complete complex tasks by providing it with available skills, workflows, and tools.
One important research question is how to improve the harness continuously.
Specifically, after an agent is deployed, we aim to continually evolve its harness by learning from past experiences, which in turn improves its performance on future tasks.

Prior work has proposed various methods for evolving the agent harness \citep{zhou2022ape,yang2023opro,khattab2023dspy,yuksekgonul2024textgrad,agrawal2025gepa,hu2024adas,lee2026metaharness}.
However, these methods rely on scoring against a validation set to guide the improvements.
In practical deployment scenarios, it is often difficult to collect a validation set that accurately estimates the distribution of future tasks to validate the updated harness.
On the other hand, the continuous operation of an agent naturally produces a rich set of trajectories from past tasks.
This leads to our central question. \textit{Can we improve the agent harness to enhance future performance when we only have access to past trajectories?}

To address this problem, we propose \textit{Retrospective Harness Optimization} (\method{}), a self-supervised method that optimizes the harness through a retrospective analysis of past trajectories.
This method employs the agent's internal self-preference over trajectories to guide the optimization process.
Figure~\ref{fig:rho-comparison} contrasts \method{} with conventional validation-feedback optimization, which iterates against a labeled validation set.
\begin{figure*}[!t]
\centering
\includegraphics[width=\textwidth]{figures/fig2-pipeline.pdf}
\caption{\looseness=-1 The \method{} pipeline. \emph{Coreset Selection} picks a small, difficulty-diverse subset of past tasks with a determinantal point process (DPP). \emph{Group Rollout} re-solves each task $G$ times and diagnoses within-trajectory failures (self-validation) and cross-trajectory disagreements (self-consistency). \emph{Harness Proposal} samples $N$ candidate harnesses and keeps the one whose rollouts are most preferred over the baseline. No ground-truth labels are used.}
\vspace{-1em}
\label{fig:rho-pipeline}
\end{figure*}
\looseness=-1 Figure~\ref{fig:rho-pipeline} illustrates this process.
Specifically, given a large set of past trajectories, \method{} first selects a diverse and challenging coreset of tasks.
Then the agent re-attempts each task in the coreset multiple times to generate parallel trajectories.
Building on this, we extract two diagnostic signals, namely self-validation within a trajectory and self-consistency across parallel trajectories.
These signals are then used to instruct the generation of harness updates.
Finally, by using the agent's pairwise self-preference, we select the most promising harness from the newly generated proposals.

\looseness=-1 We evaluate the effectiveness of \method{} across three agent domains that span software engineering, technical work, and knowledge work.
\method{} consistently improves the agent's performance across all three domains.
Notably, by running a single round of retrospective harness optimization on software-engineering trajectories, we improve the pass rate on SWE-Bench Pro \citep{deng2025swebenchpro} from 59\% to 78\%, without depending on grading against a validation set.

\looseness=-1 Furthermore, we provide a detailed analysis on how the retrospective optimization process improves performance.
We observe that \method{} designs specific skills and tools targeting typical failure modes encountered in past tasks. These components reshape the agent's action patterns and help it sustain higher accuracy in long-horizon sessions.
Additionally, we quantitatively analyze the contributions of the diagnostic signals during the retrospective process.
This analysis demonstrates that each step in \method{} progressively isolates signals that contribute to performance improvements.

\looseness=-1 Our contributions are as follows:
\begin{itemize}[label=$\diamond$]
    \item We propose retrospective harness optimization, which addresses the gap of improving the full harness (including memory, context, skills, and tools) exclusively from unlabeled trajectories.
    \item We evaluate \method{} across three scenarios and show that retrospective analysis consistently outperforms straightforward experience accumulation and reaches quality comparable to validation-feedback-driven evolution with no labels at roughly a third of its compute budget.
    \item We provide a quantitative analysis of the impact of harness optimization on agent performance, showing that gathering effective improvement signals leads to targeted changes in the harness and optimizes the agent's behavior.
\end{itemize}
